% Generated by scripts/make_paper_tables.py from the frozen results.
% Do not edit by hand; the next build overwrites it.
\begin{table}[tbp]
\centering
\caption{Repricing error against quoted bid-ask width, by surface, in implied volatility points}
\label{tab:validation}
\small
\begin{tabular}{l r r r r r}
\toprule
Surface & Options & Inside spread (\%) & Median error & Median spread & Error/spread \\
\midrule
DJX & 518 & 96.3 & 0.067 & 0.928 & 0.07 \\
NDX & 1,074 & 98.6 & 0.111 & 0.561 & 0.20 \\
RUT & 845 & 96.0 & 0.097 & 0.368 & 0.26 \\
SMH & 921 & 99.5 & 0.173 & 2.157 & 0.08 \\
SPX & 2,378 & 46.3 & 0.131 & 0.079 & 1.67 \\
\bottomrule
\end{tabular}
\notetext{\textit{Notes:} Every option used in the calibration is repriced through the local-volatility engine and the result converted back to implied volatility, so the error is in the units the market quotes. \textit{Inside spread} is the share of repriced options falling within the quoted bid-ask. The final column is the median error divided by the median quoted bid-ask width. Error is measured from the mid and the quoted band runs half a width either side of it, so a ratio below one half puts the model inside the band the market itself prices in. Four surfaces return 0.26 or better; the SPX surface returns 1.67, on a quoted width 4.7 times narrower than the next tightest here, which \S\ref{sec:validation} attributes to that market's liquidity rather than to the fit. Source: \texttt{src/vol/reprice\_check.py}.}
\end{table}
